% ============================================================
% 第9章补充：经典例题与自学元素
% 插入位置：第9章现有miniquiz之后（\end{miniquiz} 之后）
% ============================================================

\section*{习题精选与详解}
\addcontentsline{toc}{section}{习题精选}

\begin{example}{瑞利--泰勒不稳定性的增长率估算}{rayleigh-taylor-calc}
在磁约束聚变装置的"坏曲率"区域，等效重力加速度约为 $g \approx 10^{5}\,\mathrm{m/s^2}$（由曲率漂移产生），密度梯度对应的Atwood数 $A = 0.3$。界面处波长为 $\lambda = 2\,\mathrm{cm}$ 的扰动被激发。求：
\begin{enumerate}[label=(\arabic*)]
\item 该扰动的增长率 $\gamma$；
\item 振幅\textbf{增长 10 倍}所需的时间，以及振幅增长 $\ee^{10}$ 倍（即 $10$ 个 e-folding）所需的时间；
\item 若波长缩短为 $\lambda = 0.5\,\mathrm{cm}$，增长率如何变化？
\end{enumerate}
\tcblower
\textbf{解答：}

瑞利--泰勒不稳定性增长率公式为：
\[
\gamma = \sqrt{A k g}
\]
其中 $k = 2\pi/\lambda$ 为波数，$A = (\rho_2 - \rho_1)/(\rho_2 + \rho_1)$ 为Atwood数。

\textbf{(1) 计算增长率：}

波数：
\[
k = \frac{2\pi}{\lambda} = \frac{2\pi}{0.02} = 314\,\mathrm{m}^{-1}
\]

代入公式：
\[
\gamma = \sqrt{0.3\times 314\times 10^{5}} = \sqrt{9.42\times 10^{6}}
\approx 3.07\times 10^{3}\,\mathrm{s}^{-1}
\]

\textbf{(2) 增长 10 倍 vs 10 个 e-folding——两个不同的目标：}

振幅增长满足 $A(t) = A_0\ee^{\gamma t}$。

\textbf{目标一：振幅增长 10 倍}（即 $\ee^{\gamma t}=10$）：
\[
t_{10\times} = \frac{\ln 10}{\gamma} = \frac{2.303}{3.07\times 10^{3}} \approx 7.5\times 10^{-4}\,\mathrm{s} = 0.75\,\mathrm{ms}
\]

\textbf{目标二：增长 $\ee^{10}$ 倍}（即 $10$ 个 e-folding，$\gamma t=10$）：
\[
t_{10e} = \frac{10}{\gamma} = \frac{10}{3.07\times 10^{3}} \approx 3.3\times 10^{-3}\,\mathrm{s} = 3.3\,\mathrm{ms}
\]

两者相差 $\ln10\approx2.3$ 倍，\textbf{不能互换}。一个 e-folding 的时间为 $\tau=1/\gamma\approx0.33\,\mathrm{ms}$。

\textbf{约定提醒：} 本书全程用 $\exp[\ii(\vect{k}\cdot\vect{r}-\omega t)]$，因此 $\gamma=\Im\omega>0$ 对应振幅 $\propto\ee^{\gamma t}$\emph{增长}。凡涉及"增长多少倍""多少个 e-folding"，含义必须写明，因为十倍与十个 e-folding 是不同的事。

\textbf{(3) 波长缩短为 $\lambda = 0.5\,\mathrm{cm}$：}

波数变为：
\[
k' = \frac{2\pi}{0.005} = 1257\,\mathrm{m}^{-1}
\]

新的增长率：
\[
\gamma' = \sqrt{0.3\times 1257\times 10^{5}} = \sqrt{3.77\times 10^{7}}
\approx 6.14\times 10^{3}\,\mathrm{s}^{-1}
\]

增长率变为原来的 $2$ 倍（$\gamma' = \sqrt{4}\,\gamma = 2\gamma$）。

\textbf{物理分析：}
\begin{itemize}[leftmargin=*,itemsep=0pt]
\item 瑞利--泰勒增长率与 $\sqrt{k}$ 成正比，短波扰动增长更快。
\item 但极短波扰动（$k$ 很大）可能受到其他效应稳定（如表面张力、粘性、或磁场沿界面的张力）。在理想MHD中，没有这些截断机制，短波会无限增长。
\item 在托卡马克中，气球模（瑞利--泰勒型）在"坏曲率"区（外中平面）发展最快，是限制压强比 $\beta$ 的主要因素之一。
\item 约 $0.75\,\mathrm{ms}$ 内振幅增长 $10$ 倍（$3.3\,\mathrm{ms}$ 内增长 $\ee^{10}$ 倍），说明该不稳定性非常剧烈，若不加以控制（如通过理想导体壁或优化位形），等离子体会迅速损失。
\end{itemize}
\end{example}

\begin{example}{双流不稳定性的色散关系与不稳定性条件}{two-stream-calc}
\textbf{模型声明（先读）：} 本题研究\textbf{对称双流模型}——两束等密度的冷电子以 $\pm u/2$ 相向运动，净电荷中和由静止离子背景提供（离子只作电中性背景，不参与振荡）。相对漂移速度为 $u = 2\times 10^{7}\,\mathrm{m/s}$。设\textbf{总}电子密度对应的等离子体频率为 $\omega_{pe} = 5\times 10^{10}\,\mathrm{rad/s}$（每束为 $\omega_{pe}^2/2$）。求：(1) 色散关系；(2) 不稳定波数范围；(3) 最大增长率及对应波数。

\textbf{注意：} 这与"电子束穿过可动离子"（Buneman 模型，见正文 \S9.3 模型 B）是\textbf{两个不同的模型}，不可混用；本题只用对称双流模型 A。离子平衡速度为零\emph{不等于}离子扰动响应为零，但本模型里确实令离子只作背景——这个假设必须写明，不能悄悄换。
\tcblower
\textbf{解答：}

\textbf{(1) 色散关系：} 冷等离子体（$T_e=T_i=0$），每束贡献 $\omega_{pe}^2/2$：
\[
1=\frac{\omega_{pe}^2/2}{(\omega-ku/2)^2}+\frac{\omega_{pe}^2/2}{(\omega+ku/2)^2}
\]

\textbf{化为标准形式。} 令 $\Omega=\omega-ku/2$，$a=\dfrac{ku}{2\omega_{pe}}$。代入并整理可得
\[
\frac{\Omega^2}{\omega_{pe}^2}=a^2+1\pm\sqrt{1+4a^2}
\]
只有一个根给出 $\Omega^2<0$，即复频率。\textbf{不稳定要求 $\Omega^2<0$}，它发生在
\[
0<|ku|<2\sqrt2\,\omega_{pe}
\]

\textbf{(2) 不稳定波数范围：} 由 $|ku|<2\sqrt2\,\omega_{pe}$，
\[
k_{\max}=\frac{2\sqrt2\,\omega_{pe}}{u}=\frac{2\sqrt2\times 5\times10^{10}}{2\times10^{7}}
\approx 7.07\times 10^{3}\,\mathrm{m}^{-1}
\]
对应最小不稳定波长
\[
\lambda_{\min}=\frac{2\pi}{k_{\max}}=\frac{2\pi}{7.07\times10^{3}}\approx 8.9\times 10^{-4}\,\mathrm{m}\approx 0.89\,\mathrm{mm}
\]
\textbf{波长大于约 $0.89\,\mathrm{mm}$ 的模式不稳定}；更短的稳定。（常见错误是把阈值记成 $|ku|<\omega_{pe}$，对应 $k<2.5\times10^3\,\mathrm{m^{-1}}$，那是错的——复核例：取 $ku=2\omega_{pe}$ 仍有实的正增长率。）

\textbf{(3) 最大增长率：}
\[
\gamma_{\max}=\frac{\omega_{pe}}{2}=2.5\times 10^{10}\,\mathrm{rad/s}
\]
\textbf{发生在 $|ku|=\sqrt3\,\omega_{pe}$}，即
\[
k_{\text{opt}}=\frac{\sqrt3\,\omega_{pe}}{u}=\frac{1.732\times5\times10^{10}}{2\times10^{7}}\approx 4.33\times10^{3}\,\mathrm{m}^{-1}
\]
对应增长时间
\[
\tau=\frac{1}{\gamma_{\max}}\approx 4.0\times 10^{-11}\,\mathrm{s}\approx 0.040\,\mathrm{ns}=40\,\mathrm{ps}
\]
\textbf{注意 $\gamma_{\max}$ 的位置不是阈值：}阈值在 $2\sqrt2\,\omega_{pe}$，最大增长在 $\sqrt3\,\omega_{pe}$，二者不是同一个 $k$。

\textbf{物理分析：}
\begin{itemize}[leftmargin=*,itemsep=0pt]
\item 双流不稳定性将电子的定向漂移动能转化为等离子体振荡能（Langmuir波）。
\item 增长时间仅约 $40\,\mathrm{ps}$，说明在强电流等离子体中，这种不稳定性几乎是瞬发的。
\item 实际等离子体有温度：当相对漂移 $u\lesssim v_{te}$ 时，热展宽会稳定束流。
\item 束流能量转化为波，波再通过 Landau 阻尼加热电子，等效于电阻耗散——这是反常电阻的来源之一（对应冷束穿过可动离子的 Buneman 模型给出的电流驱动版本）。
\end{itemize}
\end{example}

\begin{example}{sausage模与扭曲模的定性分析}{mhd-modes}
在柱形等离子体（如Z箍缩或直圆柱 tokamak 简化模型）中，轴向电流 $I$ 产生环向磁场 $B_\theta(r)$。请分别描述 sausage模 和 扭曲模（kink模）的物理图像，并分析：
\begin{enumerate}[label=(\arabic*)]
\item 哪种模在 $m=0$ 时出现？
\item 哪种模在 $m=1$ 时出现？
\item 为什么扭曲模需要导体壁来稳定？
\end{enumerate}
\tcblower
\textbf{解答：}

\textbf{sausage模（香肠模）——$m=0$：}

 sausage模是轴对称的径向扰动（$m=0$，无角向变化），其物理图像如下：

想象一根载流等离子体柱，在某一位置受到径向压缩，柱半径略微减小。根据电流连续性 $I = \int j_z \diff S$，如果压缩是均匀的，电流密度 $j_z$ 在该处增加，环向磁场 $B_\theta = \mu_0 I(r)/(2\pi r)$ 的梯度增大。

磁压 $p_B = B_\theta^2/(2\mu_0)$ 在柱表面处增大，产生额外的向内磁压。如果等离子体压强不足以抵抗，该处会被进一步压缩——像香肠被绳子勒紧一样，形成周期性的"细腰"和"鼓胀"。

\textbf{ sausage模不稳定的条件：}
\[
\gamma^2 = \frac{B_\theta^2}{\mu_0\rho R^2}\left(2\pi\frac{R}{L}\right)^2 - \frac{2p}{\rho R^2}
\]
当磁压梯度超过压强恢复力时，不稳定性发展。 sausage模在Z箍缩中尤为危险，因为纯 $Z$ 箍缩没有纵向磁场 $B_z$ 来提供径向恢复力。

\textbf{扭曲模（kink模）——$m=1$：}

扭曲模是非轴对称的角向扰动（$m \geq 1$），其物理图像如下：

想象等离子体柱像一根柔软的绳子，被轴向电流产生的环向磁场包围。如果柱轴发生微小弯曲（$m=1$ 的螺旋形扰动），弯曲凹处的磁场线被压缩，磁压增强；凸处的磁场线被拉伸，磁压减弱。磁压差进一步将等离子体推向凹处，使弯曲加剧——正反馈！

\textbf{扭曲模的不稳定机制：}
电流产生的环向磁场在柱外随 $1/r$ 衰减。当柱弯曲时，凹侧（靠近轴）磁场更强，凸侧（远离轴）磁场更弱。磁场对等离子体柱的作用力（磁张力）在凹侧更大，将柱进一步推向凹侧，形成螺旋形扭曲。

\textbf{回答题目：}
\begin{enumerate}[label=(\arabic*)]
\item sausage模在 $m=0$ 时出现。它是轴对称的径向收缩/膨胀，无角向扭曲。
\item 扭曲模在 $m=1$ 时出现。它表现为等离子体柱轴的整体弯曲和螺旋形扭曲。
\item 导体壁稳定扭曲模的原理：理想导体壁不允许垂直于壁的磁场变化，在等离子体与壁之间的间隙中感应出镜像电流。镜像电流产生的磁场对等离子体柱施加一个恢复力，抵抗弯曲。当壁距离足够近时，这个恢复力可以超过驱动扭曲模的力，使系统稳定。这就是为什么 tokamak 中需要导体壁（或被动/主动稳定板）来抑制 $m=1$ 扭曲模。
\end{enumerate}

\textbf{对比总结：}
\begin{center}
\begin{tabular}{lcc}
\toprule
\textbf{特性} & \textbf{sausage模} ($m=0$) & \textbf{扭曲模} ($m=1$) \\
\midrule
对称性 & 轴对称 & 螺旋非对称 \\
扰动形式 & 半径周期性收缩/膨胀 & 柱轴整体弯曲 \\
稳定方法 & 纵向磁场 $B_z$（提供磁张力） & 导体壁、理想导体壳 \\
典型装置 & Z箍缩（无$B_z$时最危险） & Tokamak、RFP \\
\bottomrule
\end{tabular}
\end{center}
\end{example}

\begin{example}{漂移波不稳定性的简单估算}{drift-wave-estimate}
在密度不均匀等离子体中，密度梯度尺度长度为 $L_n = |n/(\diff n/\diff x)| = 0.1\,\mathrm{m}$，磁场 $B = 2\,\mathrm{T}$，电子温度 $T_e = 100\,\mathrm{eV}$。估算漂移波的频率和增长率。
\tcblower
\textbf{解答：}

\textbf{步骤1：计算电子热速度和漂移速度。}

电子热速度：
\[
v_{te} = \sqrt{\frac{k_{\!B}T_e}{m_e}} = \sqrt{\frac{100\times 1.602\times 10^{-19}}{9.109\times 10^{-31}}}
= \sqrt{1.76\times 10^{13}} \approx 4.19\times 10^{6}\,\mathrm{m/s}
\]

（注：这里用 $v_{te} = \sqrt{k_{\!B}T_e/m_e}$，而非 $\sqrt{2k_{\!B}T_e/m_e}$。）

电子逆磁漂移速度（$\vect{E}\times\vect{B}$ 漂移的类比，由密度梯度驱动）：
\[
v_{de} = \frac{k_{\!B}T_e}{eBL_n}
\]

将 $T_e = 100\,\mathrm{eV}$ 转换为焦耳：$k_{\!B}T_e = 1.602\times 10^{-17}\,\mathrm{J}$。
\[
v_{de} = \frac{1.602\times 10^{-17}}{1.602\times 10^{-19}\times 2\times 0.1}
= \frac{1.602\times 10^{-17}}{3.204\times 10^{-20}}
\approx 5.0\times 10^{2}\,\mathrm{m/s} = 500\,\mathrm{m/s}
\]

\textbf{步骤2：估算漂移波频率。}

漂移波频率的典型量级为：
\[
\omega_* \sim k_y v_{de}
\]

取垂直于磁场和梯度方向的波数 $k_y \sim 1/L_n = 10\,\mathrm{m}^{-1}$：
\[
\omega_* \sim 10\times 500 = 5\times 10^{3}\,\mathrm{rad/s}
\]

对应频率：
\[
f_* = \frac{\omega_*}{2\pi} \sim 800\,\mathrm{Hz}
\]

\textbf{步骤3：这一步\emph{到此为止}——不能由 $\omega_*$ 推出增长率。}

上面的 $\omega_*\sim5\times10^{3}\,\mathrm{rad/s}$ 是\textbf{实频率}（波的振荡快慢），\textbf{不}是增长率。$\omega_*=k_yv_{de}$ 只告诉你有梯度时存在一个以逆磁速度传播的模，\textbf{它本身不含任何关于稳定性的信息}。"有梯度"与"该模一定增长"之间还差一整套机制。

\textbf{要判断是否增长，必须给出非绝热的电子响应。} 漂移波不稳定的标准图景是：电子沿磁力线的响应若\emph{恰好}是绝热的（$n_{e1}/n_0=e\phi_1/k_BT_e$，即响应瞬时、无相位延迟），则密度涨落与电势涨落\emph{同相}，$\vect{E}\times\vect{B}$ 漂移只把等离子体来回搬运，净输运为零——\textbf{不增长}。只有当电子响应\emph{非绝热}、产生一个相位差 $\delta$ 时，输运才与密度涨落同相，才能把能量喂给波。产生相位差的机制包括：
\begin{itemize}[leftmargin=*,itemsep=0pt]
\item \textbf{碰撞}（耗散型漂移波）：$\delta\sim\nu_{ei}/\omega_*$，增长率 $\gamma\propto\nu_{ei}$；
\item \textbf{有限 $k_\parallel$}（朗道型/平板漂移波）：沿磁力线的波数使电子不能瞬时平衡，$\gamma\sim k_\parallel v_{te}\sqrt{\pi/2}$；
\item \textbf{曲率/坏曲率与 $\nabla T$}（ITG/TEM 等）：需要这些额外自由能源才能致不稳定。
\end{itemize}

\textbf{结论：} 本题应读作"\textbf{实频率尺度的估算}"，到此为止。$\gamma$ 必须来自具体的非绝热模型，\textbf{不能写成 $\gamma\sim\omega_*$}。请把 \S9.5 的模型假设写清楚（哪个非绝热机制？$k_\parallel$ 多大？有无碰撞？）之后再求虚部。

\textbf{物理分析：}
\begin{itemize}[leftmargin=*,itemsep=0pt]
\item 漂移波的频率远小于离子回旋频率（$\omega_{ci} = eB/m_i \approx 1.9\times 10^{8}\,\mathrm{rad/s}$ for protons），满足低频近似。
\item 漂移波\emph{可能}不稳定的核心机制：密度涨落引起电势涨落，电势涨落的 $\vect{E}\times\vect{B}$ 漂移进一步放大密度涨落——但这需要上面列出的相位差机制才能真正闭合成正反馈；无相位差时是中性振荡。
\item 在托卡马克中，漂移波及其非线性发展（漂移波湍流）是粒子反常输运的主要来源之一。
\item 虽然上述估算非常粗糙（真实色散关系涉及 $k_{\parallel}$、碰撞率、磁剪切等），但它给出了实频率的正确量级：千赫兹到兆赫兹。
\end{itemize}
\end{example}

\begin{figure}[htbp]
\centering
\begin{tikzpicture}[scale=1.0, >=Stealth, font=\small,
  box/.style={draw, rounded corners, fill=blue!10, minimum width=3.2cm, minimum height=0.9cm, align=center}]
    % 四个环节的正反馈回路
    \node[box, fill=orange!15] (dn) at (0,2.2) {密度涨落 $\delta n$};
    \node[box, fill=blue!15] (dp) at (5.6,2.2) {电势涨落 $\delta\phi$};
    \node[box, fill=green!15] (exb) at (5.6,-0.6) {$\vect{E}\times\vect{B}$ 漂移\\（扰动对流）};
    \node[box, fill=red!15] (amp) at (0,-0.6) {密度涨落被放大};

    \draw[->, very thick] (dn) -- node[above] {沿 $y$ 方向} (dp);
    \draw[->, very thick] (dp) -- node[right] {$\delta\vect{E}\perp\vect{B}$} (exb);
    \draw[->, very thick] (exb) -- node[below left=0.1cm and -0.5cm, align=center, font=\footnotesize] {把高密度区\\输运到低密度区} (amp);
    \draw[->, very thick, dashed] (amp) -- node[left] {正反馈} (dn);

    % 相位差（驱动增长的必需条件）
    \node[box, fill=yellow!25, minimum width=4.6cm] (ph) at (2.8,0.9) {\footnotesize 驱动增长必需：\\ $\delta n$ 与 $\delta\phi$ 之间的\textbf{相位差} $\delta\neq0$\\（非绝热电子响应：碰撞 / 有限 $k_\parallel$）};
    \draw[->, dashed, thick, orange!80!black] (ph.west) -- (amp.east);
    \draw[->, dashed, thick, orange!80!black] (ph.east) -- (exb.west);

    % 梯度与磁场示意（放在回路外侧，不与连线重叠）
    \node[font=\small, align=left] at (0,0.85) {};
    \node[font=\small, align=center] at (2.8,3.3) {背景：$\nabla n$ 沿 $-x$ 方向，$\vect{B}$ 垂直纸面向外};
\end{tikzpicture}
\caption{漂移波（可能）不稳定性的正反馈回路。密度梯度（沿 $-x$ 方向）中的密度涨落产生电势涨落，电势涨落的 $\vect{E}\times\vect{B}$ 漂移把高密度等离子体输运到低密度区，进一步放大密度涨落。\textbf{但回路闭合还差一环}：只有 $\delta n$ 与 $\delta\phi$ 之间存在相位差 $\delta\neq0$（即电子响应非绝热，由碰撞或有限 $k_\parallel$ 提供）时，输运才与密度涨落同相、才能把能量喂给波；若相位差为零，$\vect{E}\times\vect{B}$ 漂移只来回搬运、净输运为零，是中性振荡而非增长。磁场 $\vect{B}$ 垂直纸面向外。}
\label{fig:drift-wave-feedback}
\end{figure}

\begin{figure}[htbp]
\centering
\begin{tikzpicture}[
  scale=0.95, font=\small,
  level 1/.style={sibling distance=6cm, level distance=2cm},
  level 2/.style={sibling distance=3cm, level distance=2cm},
  level 3/.style={sibling distance=2.5cm, level distance=1.8cm},
  every node/.style={align=center, draw, rounded corners, fill=blue!10, minimum width=2cm, minimum height=0.8cm}
]
\node[fill=blue!20, minimum width=3cm] {等离子体不稳定性}
  child { node[fill=red!15] {宏观\\MHD不稳定性}
    child { node[fill=red!10] {理想MHD}
      child { node[fill=red!5] {sausage模} }
      child { node[fill=red!5] {kink\\扭曲模} }
      child { node[fill=red!5] {气球模} }
    }
    child { node[fill=orange!10] {电阻\\不稳定性}
      child { node[fill=orange!5] {撕裂模} }
      child { node[fill=orange!5] {磁重联} }
    }
  }
  child { node[fill=green!15] {微观\\动理学不稳定性}
    child { node[fill=green!10] {速度空间}
      child { node[fill=green!5] {bump-on-tail} }
      child { node[fill=green!5] {loss-cone} }
    }
    child { node[fill=teal!10] {漂移波}
      child { node[fill=teal!5] {ITG} }
      child { node[fill=teal!5] {ETG} }
      child { node[fill=teal!5] {TEM} }
    }
  };
\end{tikzpicture}
\caption{等离子体不稳定性分类树状图。宏观MHD不稳定性由压强和磁场的整体配置驱动；微观动理学不稳定性由速度空间或位形空间分布的非热力学平衡驱动。ITG = 离子温度梯度模，ETG = 电子温度梯度模，TEM = 捕获电子模。}
\end{figure}

\begin{figure}[htbp]
\centering
\begin{tikzpicture}[scale=1.1,font=\small]
% 背景盒子：重流体在上
\fill[blue!30] (-4,2) rectangle (4,4);
\node[blue!50!black] at (0,3) {\large 重流体 ($\rho_2$)};

% 背景盒子：轻流体在下
\fill[orange!30] (-4,-2) rectangle (4,2);
\node[orange!50!black] at (0,0) {\large 轻流体 ($\rho_1$)};

% 界面
\draw[thick] (-4,2) -- (-2.5,2);
\draw[thick,domain=-2.5:2.5,samples=50] plot({\x},{2+0.5*sin(2*\x r)});
\draw[thick] (2.5,2) -- (4,2);

% 重力
\draw[->,very thick,red] (0,3.5) -- (0,2.5);
\node[red,right] at (0.1,3) {$\vect{g}$};

% 扰动增长箭头
\draw[->,thick,green!50!black] (-1.5,2.5) -- (-1.5,3.2);
\draw[->,thick,green!50!black] (1.5,1.5) -- (1.5,0.8);
\node[green!50!black,right] at (1.7,1.2) {\footnotesize 扰动增长};

% 尖峰处标注
\node[align=center] at (-1.5,1.5) {\footnotesize 尖峰处重流体\\\footnotesize 下沉};
\node[align=center] at (1.5,2.5) {\footnotesize 谷处轻流体\\\footnotesize 上升};

% 波数标注
\draw[<->,thick] (-1.57,4.3) -- (1.57,4.3);
\node[above] at (0,4.3) {$\lambda = 2\pi/k$};

% 说明框
\node[draw,fill=white,align=left,rounded corners] at (0,-1.2) {
\footnotesize 瑞利--泰勒判据：$\gamma = \sqrt{A k g}$\\
\footnotesize Atwood数：$A = (\rho_2 - \rho_1)/(\rho_2 + \rho_1)$
};
\end{tikzpicture}
\caption{瑞利--泰勒不稳定性物理图像。重流体在轻流体上方，在有效重力 $g$ 作用下，界面扰动（尖峰和谷）被放大。尖峰处重流体下沉，谷处轻流体上升，形成反馈循环，扰动指数增长。}
\end{figure}

% ch9 新例题：交换不稳定性的物理判断
\begin{example}{磁流体气球模的物理图像}{ballooning-concept}
托卡马克中压强梯度驱动的气球模（ballooning mode）为什么在外中平面（坏曲率区）发展最快？用磁流体能量原理定性说明，并解释为什么提高比压 $\beta$ 会触发气球模。
\tcblower
\textbf{解答：}

\textbf{坏曲率与好曲率：} 托卡马克磁场由环向场 $B_\varphi$ 与极向场 $B_\theta$ 合成，磁力线绕环螺旋前进。环内侧（内中平面）磁场强、磁力线曲率指向等离子体内部（\textbf{好曲率}）——磁张力对扰动起稳定作用。环外侧（外中平面）磁场弱、曲率离心向外（\textbf{坏曲率}）——扰动"鼓包"向外鼓时，磁张力反而做正功。

\textbf{能量原理定性分析：} 理想 MHD 能量原理把扰动能量写成 $\delta W = \delta W_{\text{磁张力}} + \delta W_{\text{压强}}$。在坏曲率区：
\begin{itemize}[leftmargin=*,itemsep=0pt]
\item 磁张力项：扰动沿磁力线拉伸磁场，消耗能量（稳定贡献），但在坏曲率区被离心效应部分抵消；
\item 压强项：扰动把高压等离子体推到低磁压区（外中平面），释放膨胀能（不稳定贡献 $\propto \nabla p\cdot\xi$）。
\end{itemize}
当压强梯度足够大时，坏曲率区压强项的不稳定贡献超过磁张力项的稳定贡献，$\delta W<0$，气球模增长——扰动像"气球"一样从外中平面鼓出，故得此名。

\textbf{与比压 $\beta$ 的关系：} 压强项正比于 $\beta = 2\mu_0 p/B^2$（热压与磁压之比）。$\beta$ 升高意味着热压相对磁压增大，压强驱动增强而磁场稳定能力相对减弱。当 $\beta$ 超过临界值 $\beta_c$（工程上通常 $\beta_c\sim3\%$ 量级，取决于位形与剪切断面）时，气球模被触发——这正是托卡马克比压极限（Troyon 极限）的物理根源：$\beta$ 不能无限提高，否则压强梯度驱动的气球模立即破坏约束。

\textbf{工程对策：} 提高磁场剪切断面（优化极向场位形）、降低边缘压强梯度（H模台地）、或使用导体壁靠近等离子体（ resistive wall 稳定低频模）。现代托卡马克通过这些手段把 $\beta$ 推到 $\sim3\%--5\%$，接近磁约束聚变堆的经济运行区间。
\end{example}

\begin{formulabox}{第9章核心公式速查}{ch9-formulas}
\begin{itemize}[leftmargin=*,itemsep=0pt]
\item 不稳定性判据：$\text{Im}(\omega) > 0$，增长率 $\gamma = \text{Im}(\omega)$
\item 瑞利--泰勒增长率：$\displaystyle\gamma = \sqrt{A k g}$，$A = \frac{\rho_2 - \rho_1}{\rho_2 + \rho_1}$
\item 对称双流（两束等密度冷电子，相对漂移 $u$）：$\displaystyle 1=\frac{\omega_{pe}^2/2}{(\omega-ku/2)^2}+\frac{\omega_{pe}^2/2}{(\omega+ku/2)^2}$，即 $\dfrac{\Omega^2}{\omega_{pe}^2}=a^2+1\pm\sqrt{1+4a^2}$（$\Omega=\omega-ku/2$，$a=ku/2\omega_{pe}$）
\item 对称双流不稳定条件：$\displaystyle 0<|ku|<2\sqrt2\,\omega_{pe}$（\textbf{不是} $ku<\omega_{pe}$）
\item 对称双流最大增长率：$\displaystyle\gamma_{\max}=\frac{\omega_{pe}}{2}$，发生在 $\displaystyle|ku|=\sqrt3\,\omega_{pe}$（\textbf{不是阈值处}）
\item Buneman（电子束穿过可动离子）：$\displaystyle 1=\frac{\omega_{pe}^2}{(\omega-ku_e)^2}+\frac{\omega_{pi}^2}{(\omega-ku_i)^2}$，与对称双流是两个模型，不可混用
\item  sausage模（$m=0$）：轴对称径向收缩，纵向磁场 $B_z$ 提供稳定
\item 扭曲模（$m=1$）：柱轴螺旋弯曲，导体壁镜像电流提供稳定
\item 漂移波频率（估算）：$\displaystyle\omega_* \sim k_y v_{de} = k_y \frac{k_{\!B}T_e}{eBL_n}$
\item 有效重力（曲率漂移）：$\displaystyle g \sim \frac{v_{th}^2}{R_c}$
\item 反常输运：湍流电场导致随机散射，输运系数远大于经典值
\end{itemize}
\end{formulabox}

\begin{checklistbox}{第9章自学检查清单}{ch9-checklist}
\begin{itemize}[leftmargin=*,label=$\square$]
\item 我能区分宏观MHD不稳定性（压强/磁场配置驱动）和微观动理学不稳定性（速度/位形空间驱动）
\item 我能从瑞利--泰勒公式估算增长率，并解释为什么短波增长更快
\item 我能写出双流不稳定性的色散关系，并判断不稳定波数范围
\item 我能在物理图像上区分 sausage模（$m=0$）和扭曲模（$m=1$），并知道各自的稳定方法
\item 我理解漂移波的回路：密度涨落 $\to$ 电势涨落 $\to$ $\vect{E}\times\vect{B}$ 漂移 $\to$ 输运，\textbf{并且知道回路闭合还差一环}——必须有相位差 $\delta\neq0$（非绝热电子响应）才能增长；$\omega_*=k_yv_{de}$ 只是实频率，不等于增长率
\item 我能解释为什么反常输运远大于经典输运，以及它与不稳定性的关系
\item 我知道抑制不稳定性是磁约束聚变的核心挑战（如气球模限制 $\beta$、扭曲模需要壁稳定）
\end{itemize}
\end{checklistbox}

\endinput
